\section{Chebotarev density theorem}

The general Galois case is reduced to the abelian case via the cyclic
subextension generated by a class representative. Source: Sharifi 7.2.2
Step~1 (p.~143).

\begin{lemma}\label{lem:count-primes-above-frob-sigma}
  \lean{Chebotarev.count_primes_above_with_frobenius_eq_sigma}
  \uses{def:frobenius-at, def:frobenius-class}
  Let $L/K$ be a finite Galois extension with $G=\Gal{L/K}$,
  $C\subseteq G$ a conjugacy class, $\sigma\in C$, and $\fp$ a nonzero
  prime of $\OK$ unramified in $L$ with $\sigma_\fp = C$. The number
  of primes $\fP$ of $\Ocirc_L$ above $\fp$ with $\Frob_\fP = \sigma$
  is exactly $\abs{G}/(f\cdot\abs{C})$, where $f = \mathrm{ord}(\sigma)$.
\end{lemma}

\begin{proof}
  The Galois group $G$ acts transitively on primes above $\fp$, with
  stabiliser $D_\fP$. Hence the orbit has size $\abs{G}/\abs{D_\fP}$. For
  unramified $\fp$, $D_\fP=\ang{\Frob_\fP}$ has order $f$, so there are
  $\abs{G}/f$ primes above $\fp$. Their Frobenius elements form a
  $G$-orbit in $C$ (under conjugation), which is all of $C$ by
  transitivity; this orbit has size $\abs{C}$, and within each orbit
  the count of representatives $\sigma$ is $\abs{G}/(f\cdot\abs{C})$
  (Sharifi p.~143).
\end{proof}

\begin{theorem}[Chebotarev density theorem]\label{thm:chebotarev-density}
  \lean{Chebotarev.chebotarev_density}
  \uses{def:dirichlet-density, def:frobenius-class, thm:chebotarev-abelian, lem:count-primes-above-frob-sigma}
  For a finite Galois extension $L/K$ of number fields with
  $G=\Gal{L/K}$ and any conjugacy class $C\subseteq G$,
  \[
    \delta\bigl(\setof{\fp\subset\OK}{\sigma_\fp = C}\bigr)
    \;=\; \frac{\abs{C}}{\abs{G}}.
  \]
\end{theorem}

\begin{proof}
  Pick $\sigma\in C$, let $E = L^{\ang\sigma}$, and set
  $f = \mathrm{ord}(\sigma) = \abs{\ang\sigma}$. Then $L/E$ is cyclic
  (hence abelian) of degree $f$, so
  \cref{thm:chebotarev-abelian} applied to $L/E$ gives
  $\delta_E(\setof{P}{\sigma_P = \sigma}) = 1/f$.

  Combine the above with the counting from
  \cref{lem:count-primes-above-frob-sigma}: for primes $P$ of $E$ above
  $\fp \in S = \setof{\fp}{\sigma_\fp = C}$, the count with
  $\sigma_P = \sigma$ is exactly $\abs{G}/(f\cdot\abs{C})$ per orbit.
  Using $\sum_\fp N\fp^{-s} \sim \sum_P NP^{-s}$
  (\cref{lem:prime-ideal-sum-log} applied to both $K$ and $E$),
  \[
    \delta_K(S) \;=\;
    \frac{f\cdot\abs{C}}{\abs{G}}\cdot\delta_E(\setof{P}{\sigma_P = \sigma})
    \;=\;
    \frac{f\cdot\abs{C}}{\abs{G}}\cdot\frac{1}{f}
    \;=\;
    \frac{\abs{C}}{\abs{G}}.
  \]
\end{proof}

\subsection*{Corollaries}

\begin{corollary}\label{cor:density-split-completely}
  \lean{Chebotarev.density_split_completely}
  \uses{thm:chebotarev-density}
  For a finite Galois extension $L/K$ of number fields, the Dirichlet
  density of primes of $\OK$ that split completely in $L$ equals
  $1/[L:K]$.
\end{corollary}

\begin{proof}
  A prime $\fp$ splits completely in $L$ iff every $\fP \mid \fp$ has
  $\Frob_\fP = 1$. The conjugacy class $\{1\}$ has size $1$, and
  $\abs{G}=[L:K]$, so \cref{thm:chebotarev-density} gives
  $\delta = 1/[L:K]$.
\end{proof}

\begin{corollary}[Dirichlet, density form]\label{cor:dirichlet-ap-density}
  \lean{Chebotarev.dirichlet_primes_in_AP}
  \uses{thm:chebotarev-density}
  For $n \ge 1$ and $a \in (\Z/n\Z)^\times$, the Dirichlet density of
  the rational primes $p$ with $p \equiv a \pmod{n}$ equals
  $1/\varphi(n)$.
\end{corollary}

\begin{proof}
  Specialise \cref{thm:chebotarev-density} to $K = \Q$, $L = \Q(\zeta_n)$.
  The cyclotomic character identifies $\Gal{\Q(\zeta_n)/\Q}$ with
  $(\Z/n\Z)^\times$; under this identification, the Frobenius
  $\sigma_p$ at unramified $p$ corresponds to $p \bmod n$. The Galois
  group is abelian, so each conjugacy class is a singleton; the set of
  primes with $p \equiv a \pmod{n}$ is the Frobenius preimage of
  $\{a\}$, whose density is $1/\varphi(n)$.
\end{proof}
